%%%
%
% $Autor: Malena Wilken $
% $Date: 2026-01-20 $
% $File: FlowchartDeployment.tex$
% $Version: 1 $
% !TeX spellcheck = en_US
% !TeX encoding = utf8
% !TeX root = filename 
% !TeX TXS-program:bibliography = txs:///biber
%
%%%%%%%%%%%%

%source: https://projecthub.arduino.cc/paulsb/temp-and-humidity-monitor-with-graphs-and-battery-monitor-cd011a

% https://www.az-delivery.de/products/az-delivery-laderegler-tp4056-mini-usb?variant=12239811084384
%https://www.az-delivery.de/products/cd60l-batterieladecontroller
%https://www.az-delivery.de/products/mini-solarpanel?variant=39475872792672


    
\chapter{Powering the TinyML Shield}

After assembling the TinyML hardware and verifying correct operation through initial software tests, the system is typically powered directly via a USB connection to a host computer. In development state the microcontroller has been powered via the USB micro-B interface of the Nano 33 BLE Sense, where the supplied 5 V is internally regulated down to 3.3 V as the logical supply voltage of the MCU. While this approach is well suited for development and debugging, a functional prototype should not rely on continuous power delivery from a host computer. A more practical alternative is the use of an external AC–DC converter with USB output. In addition to enabling standalone operation, such converters typically provide a higher available current on the 5 V rail than a standard USB port, which is commonly limited to approximately 500 mA. This increased current capacity can be relevant when the system is extended with power-demanding components, such as audio actuators or external peripherals are often intended to operate independently of a workstation. For this reason, a  understanding of the available power supply options is essential. The following section therefore focuses on how to power the Tiny Machine Learning Shield in standalone scenarios \cite{ICTP:2022}.


\section{Battery}

While the previously described approach enables operation without a host computer, the system remains dependent on a fixed power outlet. Fully mobile operation therefore requires a battery-based power supply. The same principle used for USB operation applies in this case, as the onboard voltage regulator, specifically the MPM3610, is responsible for converting the input voltage to a stable 3.3 V supply for the microcontroller. According to its datasheet, the regulator supports input voltages ranging from 4.5 V to 21 V. The 5 V provided via USB lies within this range, and any battery-based solution must satisfy the same requirement. As a result, single-cell 3.7 V lithium batteries are not suitable without additional circuitry, whereas a 9 V alkaline battery represents a straightforward and compatible choice. Electrical connection of an external battery is facilitated by a dedicated terminal block located on the Tiny Machine Learning Shield. This connector is labeled with “VIN” and “GND,” where GND serves as the reference potential and must be connected to the negative terminal of the battery. The positive terminal is connected to VIN, which accepts input voltages within the specified range of the regulator. Using a standard 9 V battery clip with pre-stripped leads allows for a simple and reliable connection, enabling standalone operation without further modifications.



\begin{center}
	
	
\begin{tikzpicture}
		
		\coordinate (BoardOrigin) at (6.8,28.1);
		\coordinate (BoardOriginNano) at (3,43.4);
		
		\draw[black, line width=2pt] (4.4,46.21) -- (5.4,46.21) -- (5.4,44.3) -- (7,44.3);
		\draw[red, line width=2pt] (4.4,46.11) -- (5.1,46.11) -- (5.1,43.95) -- (7,43.95);
		
		\begin{scope}[scale=0.65, shift={(BoardOrigin)}]
			\ArduinoNanoShieldTikz
		\end{scope}
		
		\begin{scope}[scale=0.75, shift={(BoardOriginNano)}]
			\Battery
		\end{scope}
		
		
	\end{tikzpicture}    
	
	\captionof{figure}[Installation of a Battery with the Tiny Machine Learning Shield]{Installation of a Battery with the Tiny Machine Learning Shield}
	\label{fig:BatteryShield}
\end{center}


To connect a battery to the Tiny Machine Learning Shield, as shown in figure \ref{fig:BatteryShield}, the negative terminal of the battery must first be connected to the GND pin of the onboard terminal block, typically using a small flat-head screwdriver to secure the wire. The positive terminal is then connected to the VIN pin, completing the power connection and enabling battery-powered operation of the system. For safe and reliable use, it is recommended to disconnect the Nano 33 BLE Sense from any USB power source before attaching a battery, in order to avoid competing power inputs. When using a standard 9 V alkaline battery with a nominal capacity of approximately 550 mAh, the system can be supplied with a current of around 15 mA for an estimated duration of 37 hours, depending on the actual load and operating conditions.










\section{Checking the Battery Voltage}

An analog input pin is used to read the voltage. As we are running from a 3.7V volt battery, we need to adjust the reference voltage used by the pin as otherwise it would be comparing the voltage to itself. The statement \PYTHON{analogReference(INTERNAL)} sets the pin to compare the input voltage to a regulated 1.1V. We therefore need to reduce the voltage on the input pin to less than 1.1V for this to work. This is done by dividing the voltage using 2 resistors, 1m and 330k ohms. This divides the voltage by approximately 4 so when the battery is fully charged, which is 4.2V, the voltage at the pin input is 4.2/4 = 1.05V. 




\begin{center}
	\captionof{code}{Code to Check Battery Level}
	\lstinputlisting[style=all]{../../Code/TinyMLKit/TestBattery/TestBattery.ino}
	\label{lst:Battery}
\end{center}


The function \PYTHON{BatterVoltage()} reads the analog pin, which will range from 0 for 0V to 1,023 for 1.1V and using this reading calculates the actual voltage coming form the battery. 

The function \PYTHON{DrawScreenSave()} calls this then selects the appropriate bitmap to display based on the following: 

\begin{itemize}
    \item If voltage is greater then 3.6V - full 
    \item Voltage between 3.5 and 3.6V - 3/4 
    \item Voltage between 3.4 and 3.5V - half 
    \item Voltage between 3.3 and 3.4V - 1/4 
    \item Voltage < 3.3V - empty 
\end{itemize}



%\section{Batterieclip}
%
%Der Batterieclip in Abb. \ref{Batterieclip für 9-Volt-Block}, der vom Hersteller \textit{reichelt} ist, kann vertikal an einen 9-Volt-Block angeschlossen werden. Die dazugehörigen Anschlussdrähte haben eine Länge von 150 mm. Der Anschlussclip ist in der I-Form ausgeführt, weshalb er sich platzsparend ins Gehäuse einbinden lässt \cite{Reichelt:2011}.
%
%\begin{figure}[h]
%    \begin{center}
%        \includegraphics[width=3in]{Battery/clip.jpg}
%        \caption{Batterieclip für 9-Volt-Block\cite{Reichelt:2024a}}
%        \label{Batterieclip für 9-Volt-Block}
%    \end{center}
%\end{figure} 
%
%
%
%\section{Spannungssensor}
%
%Der Spannungssensor in Abb. \ref{Spannungssensor} von dem Hersteller \textit{Shenzhen Global Technology Co., Ltd} kann bei der Versorgungsspannung von 3,3 V Spannungen in dem Bereich von 0 V bis 16,5 V messen. Dieser wird genutzt, um den Ladestand der Batterie zu überwachen. Die analoge Auflösung des Sensors liegt bei 10 Bit. Damit kann bei dem angegebenen Messbereich die Spannung mit der Auflösung von 0,016 V gemessen werden. Zur Eingangsschnittstelle gehört der \ac{vcc}-Anschluss und der \ac{gnd}-Anschluss \cite{Shenzhen:2015}. Die Bauteilmaße betragen 13 mm x 27 mm.
%
%
%
%{\begin{center}
%		
%  \includegraphics[width=0.6\textwidth]{Battery/VoltageSensor}
%  
%  \captionof{figure}{Voltage Sensor 170640  von dem Hersteller \textit{Shenzhen Global Technology Co., Ltd}\cite{Shenzhen:2015}}   		
%		
%	\end{center}
%}
%	
%Mit dem Sketch \FILE{TestBattery.ino}, siehe \ref~{TestBattery}, soll der Spannungssensor getestet werden.
%
%\medskip
%
%
%
%\begin{code}
%	
%	\caption{Beispiel-Sketch zum Testen der Batterie}\label{TestBattery}
%
%    \ArduinoExternal{}{../../Code/Arduino/Battery/TestBattery.ino}
%\end{code}
%
%\subsection{Durchführung}
%
%Für den Test werden die folgenden Hardware-Komponenten benötigt:
%
%\begin{itemize}
%    \item Arduino Nano 33 BLE Sense Lite
%    \item Tiny Machine Learning Shield
%    \item USB-A auf USB-Mikro Verbindungskabel
%    \item Grove Jumper zu Grove 4 Pin Kabel
%    \item Spannungssensor
%    \item Batterie
%    \item Batterieclip
%\end{itemize}
%
%Die Hardware-Komponenten werden  zusammengebaut, aber die Batterie wird noch nicht angeschlossen. Dann wird der Arduino Nano 33 BLE Sense Lite mit einem Computer verbunden. Anschließend wird der Sketch \FILE{TestBattery.ino} auf den Arduino Nano 33 BLE Sense Lite geladen und der serielle Monitor in der Arduino \acs{ide} geöffnet. Die Batterie wird während des Tests an den Batterieclip angeschlossen und der gemessene Wert bei dem seriellen Monitor ausgelesen.
%
%\subsection{Ergebnisse}
%
%Zu Beginn des Tests zeigt der serielle Monitor die Spannung 0 V an. Nach dem Anschließen der Batterie an den Spannungssensor wird die Spannung 9,6 V angezeigt.  Abbildung~\ref{BildSpannungTest} zeigt den gemessenen Spannungsverlauf. Die angezeigte Spannung liegt in dem erwarteten Bereich einer 9 V Batterie.
%
%\begin{figure}[h]
%    \begin{center}
%        \includegraphics[width=\textwidth]{Battery/BatteryTest.png}
%        \caption{Testoutput des Spannungssensors}
%        \label{BildSpannungTest}
%    \end{center}
%\end{figure}
%
